\subsection{A transport criterion with the missing hypothesis made
explicit}\label{proof:positive_structure:a-transport-criterion-with-the-missing-hypothesis-made-explicit}

Let \(Z\) be uniform on \(\{-1,1\}^N\), let \(T:\{-1,1\}^N \to\{-1,1\}^n\),
and let \(\mu =T_*U_N\). Assume:

\begin{enumerate}
\def\labelenumi{\arabic{enumi}.}
\tightlist
\item
  Each output coordinate \(T_i\) has multilinear degree at most \(k\) in
  the source bits.
\item
  Each source coordinate \(j\) has a deterministic set \(S_j\) of at
  most \(L\) outputs it can affect. Thus \(T(z)\) and \(T(z^j)\) agree
  outside \(S_j\) for every \(z\).
\item
  Choose a fixed ordering of \(S_j\), and replace these output
  coordinates one at a time to interpolate from \(T(z)\) to \(T(z^j)\).
  Every intermediate random configuration \(W_{j,l}\) has law dominated
  by \(C \mu\): for every nonnegative function \(h\),
  \(\mathbb{E} h(W_{j,l}) \le  C \mathbb{E}_\mu  h\).
\end{enumerate}

\textbf{Theorem P1 (local-map transfer).} Under these assumptions,

\[
\boxed{\max_i I_i^\mu(f)
\ge \frac{c_0}{C L^2 k^6}\frac{v^{5/2}}{d^6}.}
\tag{P1}
\]

Source coordinates with no affected outputs have zero influence and can
be ignored.

\textbf{Proof.} The composition \(F=f\circ T\), after
multilinearization, has degree at most \(kd\), is globally bounded, and
has variance \(v\). A source-bit flip is telescoped over at most \(L\)
output flips. Cauchy--Schwarz gives

\[
I_j^U(F)\le L\sum_{i\in S_j}\mathbb E(D_i f(W_{j,l(i)}))^2
\le LC\sum_{i\in S_j}I_i^\mu(f)
\le L^2C\max_i I_i^\mu(f).
\]

Apply UAA to \(F\) and rearrange. \qed

The domination hypothesis cannot be omitted. A local source flip may
move the output through configurations carrying negligible or zero
target mass, even if its endpoints have the correct marginal law.

\subsubsection{Finite conditional energies supply path
domination}\label{proof:positive_structure:finite-conditional-energies-supply-path-domination}

Suppose \(\mu\) has full support and, uniformly over all coordinates and
conditioning configurations,

\[
\min_{s\in\{-1,1\}}\mu(X_i=s\mid X_{-i}=x_{-i})\ge a>0.
\tag{FE}
\]

Then every conditional block atom on a set \(S\) of size at most \(L\)
has probability at least \(a^L\). Reveal the coordinates of \(S\)
sequentially: the partially conditioned probabilities are averages of
the fully conditioned ones and remain at least \(a\).

If \(X~\mu\) and any randomized modification \(W\) agrees with \(X\)
outside \(S\), then

\[
\Pr(W=x)\le\mu(X_{-S}=x_{-S})\le a^{-L}\mu(x).
\]

No independence between the modification and \(X\) is needed.

\textbf{Corollary P2.} Assumptions 1--2 of P1 and (FE) imply

\[
\boxed{\max_i I_i^\mu(f)
\ge \frac{c_0a^L}{L^2k^6}\frac{v^{5/2}}{d^6}.}
\tag{P2}
\]

This is dimension independent when \(a,k,L\) are fixed. A sampling
radius growing with dimension cannot silently be absorbed into the
constant.

For a classical interaction Hamiltonian

\[
H(x)=-\sum_A J_A\prod_{i\in A}x_i,
\qquad\Lambda=\max_i\sum_{A\ni i}|J_A|,
\]

the Gibbs law has (FE) with \(a=(1+\exp (2 \beta \Lambda))^{-1}\). A
bounded local-field norm therefore supplies path comparison \textbf{if
an exact sampling map with controlled \(k,L\) is also available}.
Bounded local fields alone do not construct that map.

\subsubsection{Counterexamples to tempting weaker
statements}\label{proof:positive_structure:counterexamples-to-tempting-weaker-statements}

\textbf{Local degree and overlap alone do not suffice.} Let
\(T(z)=(z,z,z)\) for one uniform source spin, and let

\[
f(x)=\operatorname{Maj}_3(x)=\frac{x_1+x_2+x_3-x_1x_2x_3}{2}.
\]

The map has output degree one and source fanout three. The globally
bounded \(f\) has degree three and variance one under \(T_*U\), but
every coordinate influence is zero: a single bit flip at \(+++\) or
\(---\) leaves majority unchanged. The interpolating paths leave the
support.

\textbf{Full support without uniform nondegeneracy does not repair the
claim.} On each independent triple take
\((1-\epsilon _n) \operatorname{Unif}\{+++ , ---\}+\epsilon _n U_3\), with
\(\epsilon _n \to  0\). Majority of the first triple still has variance
one and all its nonzero influences equal \(\epsilon _n/2\). These are
full-support laws with independent bounded blocks; a uniform lower bound
on block atom masses is essential below.

\textbf{Finite conditional energies and easy sampling alone do not
suffice.} Let \(Z\) be uniform and \(E_i\) independent spins
with mean \(t \in (0,1)\), and put \(X_i=ZE_i\). Every conditional
one-spin probability is in \([(1-t)/2,(1+t)/2]\). Nevertheless

\[
f(X)=\frac1n\sum_iX_i,\quad
\operatorname{Var}f=t^2+(1-t^2)/n,\quad I_i(f)=n^{-2}.
\]

The sampler has depth two and degree two in its independent, partly
biased source spins, but the source \(Z\) affects every output. To give
a counterexample with only fair source bits, fix \(t=1/2\) and generate
each \(E_i\) from two fair bits, assigning three source strings
to \(+1\) and one to \(-1\). Then each output has degree at most three
in fair bits, while the source \(Z\) still has fanout \(n\). With this
explicit encoding, unbounded fanout is the sole missing hypothesis in
P2.

An explicit two-body Gibbs realization includes the root as an input
spin:

\[
H(x_0,x_1,\ldots,x_n)=-Jx_0\sum_{i=1}^n x_i,
\quad t=\tanh(\beta J).
\]

Leaf magnetization gives the same counterexample, with root influence
zero, for every fixed \(\beta  J>0\). Interaction arity and individual
coupling magnitude are bounded, but the root degree and total incident
coupling are unbounded. Those first two bounds alone therefore cannot
ensure a universal positive-temperature AA regime.
